%%%%%%%%%%  scan idx 537  |  folio 524  %%%%%%%%%%
\SGAPrintPage{524}%

\[
x\cl_\bullet(\underline{O}_Z)\in\Filt_{j-i}(X).
\]
To avoid any confusion, denote by $1_Z$ the class of $\underline{O}_Z$ in $K_\bullet(Z)$. If $\alpha:Z\to X$ is the canonical immersion, then the projection formula (IV 2.11.1.2) gives
\[
x\cl_\bullet(\underline{O}_Z)=x\alpha_*(1_Z)=\alpha_*(\alpha^*(x)1_Z).
\]
One has $\alpha^*(x)\in\Filt^i(Z)$, so by (1.1.8) and (1.3.3) for $X=Z$, $i+j\leq k$, one obtains
\[
\alpha_*(\alpha^*(x)1_Z)\in\Filt_{j-i}(X),
\]
whence the assertion.

Theorem 1.3.2 is trivial for $i+j\leq0$. Suppose by induction that it has been proved for $i+j<k$ and for every noetherian scheme $X$. By the preceding remark, it remains to prove that for every \emph{integral} noetherian scheme $X$ of dimension $j$ and every product (1.2.3) satisfying (1.2.4), one has
\sgaStableTarget{sga6.e10.eq.1-3-4}\begin{equation}\tag{1.3.4}
P\cl_\bullet(\underline{O}_X)\in\Filt_{j-i}(X).
\end{equation}
If $k>1$ in (1.2.3), the induction hypothesis and the associativity of the law (1.3.1) give the result. We may therefore suppose ($X$ being connected) that
\sgaStableTarget{sga6.e10.eq.1-3-5}\begin{equation}\tag{1.3.5}
x=P=\gamma^m(\cl^{\bullet}(\underline{E})-n),
\end{equation}
where $m\geq i$ and $\underline{E}$ is a locally free $\underline{O}_X$-module of rank $n>0$.

Distinguish two cases:

I. Suppose that one can find invertible $\underline{O}_X$-modules $\underline{L}_1,\ldots,\underline{L}_n$ such that
\[
\cl^{\bullet}(\underline{E})=\cl^{\bullet}(\underline{L}_1)+\cdots+\cl^{\bullet}(\underline{L}_n)\in K^{\bullet}(X).
\]


%%%%%%%%%%  scan idx 538  |  folio 525  %%%%%%%%%%
\SGAPrintPage{525}%

It follows (V) that $x$ is the $m$-th elementary symmetric function in the $\cl^{\bullet}(\underline{L}_r)-1$ ($r=1,\ldots,n$), hence is a homogeneous polynomial of degree $m$ with integer coefficients in the variables $(\cl^{\bullet}(\underline{L}_r)-1)$ ($r=1,\ldots,n$). Applying the induction hypothesis and the associativity of the law (1.3.1) again, we are thus reduced to proving that, for every invertible $\underline{O}_X$-module $\underline{L}$, one has
\sgaStableTarget{sga6.e10.eq.1-3-6}\begin{equation}\tag{1.3.6}
(\cl^{\bullet}(\underline{L})-1)\cl_\bullet(\underline{O}_X)\in\Filt_{j-1}(X)
\end{equation}
Since the scheme $X$ is integral, $\underline{L}$ and $\underline{O}_X$ are isomorphic to submodules of the constant $\underline{O}_X$-module $\underline{R}_X$ of rational functions on $X$. The sum $\underline{O}_X$-module $\underline{F}=\underline{O}_X+\underline{L}$ is coherent, and one obtains
\begin{align*}
(\cl^{\bullet}(\underline{L})-1)\cl_\bullet(\underline{O}_X)
&=\cl_\bullet(\underline{L})-\cl_\bullet(\underline{O}_X)\\
&=(\cl_\bullet(\underline{F})-\cl_\bullet(\underline{O}_X))-(\cl_\bullet(\underline{F})-\cl_\bullet(\underline{L}))\\
&=\cl_\bullet(\underline{F}/\underline{O}_X)-\cl_\bullet(\underline{F}/\underline{L}).
\end{align*}
But the fibers at the generic point of the sheaves $\underline{F}/\underline{O}_X$ and $\underline{F}/\underline{L}$ are zero, so their supports in $X$ are nowhere dense ($X$ being irreducible), hence of dimension $\leq j-1$. Assertion (1.3.6) follows immediately.

II. It remains to prove the general case. Introduce the flag bundle
\[
g:D\longrightarrow X
\]
of $\underline{E}$ over $X$, on which one can find invertible $\underline{O}_D$-modules $\underline{L}_1,\ldots,\underline{L}_n$ such that
\sgaStableTarget{sga6.e10.eq.1-3-7}\begin{equation}\tag{1.3.7}
g^*(\cl^{\bullet}(\underline{E}))=\cl^{\bullet}(L_1)+\cdots+\cl^{\bullet}(\underline{L}_n)\in K^{\bullet}(D).
\end{equation}
Let $e$ be the generic point of $X$, and $\underline{k}(e)$ the residue field of $e$. Then one can find a rational point $u$ of the generic fiber $D_e$ of $D$


%%%%%%%%%%  scan idx 539  |  folio 526  %%%%%%%%%%
\SGAPrintPage{526}%

over $\underline{k}(e)$, since these points are in bijective correspondence with the flags of type $(1,\ldots,1)$ ($n$ terms) of the vector space $\underline{E}_e$ of dimension $n>0$ over $\underline{k}(e)$. On the other hand, these rational points also correspond to the germs of sections of $g$ in a neighborhood of $e$, so one can find such a section $s:U\to D$ over an open subset $U$ of $X$. Let $X'=\overline{s(U)}$ be the closure of $s(U)$ with its reduced induced structure, and let $f:X'\to X$ be the restriction of $g$ to $X'$. The morphism $f$ is projective and, since it is an isomorphism over $s(U)$, it is birational.
Thus $\dim(X')\leq\dim(X)$ (EGA IV 5.6.6.1). On the other hand, $f(X')=g(X')$ is closed and contains $U$, so $f$ is surjective and $\dim(X')\geq\dim(X)$ (EGA IV 5.4.1(ii)).

The hypotheses of case I are therefore satisfied for $X'$ and $x'=f^*(x)\in\Filt^i(X')$, and one obtains
\sgaStableTarget{sga6.e10.eq.1-3-7.occ2}\begin{equation}\tag{1.3.7'}
f^*(x)\cl_\bullet(\underline{O}_{X'})\in\Filt_{j-i}(X').
\end{equation}
Taking account of the projection formula and (1.1.8), one concludes that:
\sgaStableTarget{sga6.e10.eq.1-3-8}\begin{equation}\tag{1.3.8}
x f_*(\cl_\bullet(\underline{O}_{X'}))\in\Filt_{j-i}(X).
\end{equation}

On the other hand, the canonical homomorphism
\[
\underline{O}_X\longrightarrow f_*(\underline{O}_{X'})
\]
is an isomorphism over the dense open subset $U$, and the $\underline{O}_X$-modules $R^if_*(\underline{O}_{X'})$ ($i\geq1$) are concentrated on the nowhere dense closed subset $X-U$. It follows that
\sgaStableTarget{sga6.e10.eq.1-3-9}\begin{equation}\tag{1.3.9}
f_*(\cl_\bullet(\underline{O}_{X'}))=\cl_\bullet(\underline{O}_X)+w,

%%%%%%%%%%  scan idx 540  |  folio 527  %%%%%%%%%%
\end{equation}
\SGAPrintPage{527}%

where $w\in\Filt_{j-1}(X)$. By the induction hypothesis, one has $x\cdot w\in\Filt_{j-i-1}(X)$, whence, taking account of formulas (1.3.8) and (1.3.9), the assertion
\[
x\cl_\bullet(\underline{O}_X)\in\Filt_{j-i}(X),\qquad\text{Q.E.D.}
\]

\sgaStableTarget{sga6.e10.rem.1-4}\textbf{Remark 1.4.} Let $X$ be a locally noetherian scheme. Observe that, in this section, we have used \emph{only} the following properties of the ``dimension function'' $D(x)=D_Y(x)=\dim(\overline{\{x\}})$ on proper $X$-schemes $Y$:
\begin{enumerate}[label=(\roman*)]
\item For every proper $X$-scheme $Y$ and every $x\in Y$, one has $D_Y(x)\in\mathbb Z$.
\item If $y\ne x$ is a specialization of $x$, one has
\[
D(y)<D(x).
\]
\item For every $X$-morphism $f:Y_1\to Y_2$ of proper $X$-schemes and every $y\in Y_1$, one has
\[
D_{Y_2}(f(y))\leq D_{Y_1}(y).
\]
\item Equality holds in (iii) when $f$ is an immersion, or when $Y_1$ is irreducible, $f$ is birational and surjective, and $y$ is the generic point of $Y_1$.
\end{enumerate}

\uline{This means that all the results of this section, in particular} (1.3.2) \uline{and} (1.3.3)\uline{, remain valid unchanged when} ``\uline{dimension}'' \uline{is replaced by a function} $D$ \uline{satisfying conditions} (i) to (iv) \uline{above and the filtration of} $K_\bullet(X)$ \uline{is defined using such a function} $D$.


%%%%%%%%%%  scan idx 541  |  folio 528  %%%%%%%%%%
\SGAPrintPage{528}%

\sgaStableTarget{sga6.e10.ex.1-5}\textbf{Example 1.5.} Let $S$ be a locally noetherian universally catenary scheme (for example, a subscheme of a Cohen--Macaulay scheme). Put
\sgaStableTarget{sga6.e10.eq.1-5-1}\begin{equation}\tag{1.5.1}
D_S(x)=n-\codim(\overline{\{x\}},S)=n-\dimdot\underline{O}_{S,x}\qquad(x\in S),
\end{equation}
and, for every proper $S$-scheme $s_Y:Y\to S$ and $y\in Y$, also put:
\sgaStableTarget{sga6.e10.eq.1-5-2}\begin{equation}\tag{1.5.2}
D_Y(y)=D_S(s_Y(y))+\degtr_{k(s_Y(y))}k(y).
\end{equation}

If $f:Y_1\to Y_2$ is an $S$-morphism of proper $S$-schemes, the additivity of transcendence degrees and (1.5.2) immediately give
\sgaStableTarget{sga6.e10.eq.1-5-3}\begin{equation}\tag{1.5.3}
D_{Y_1}(y)=D_{Y_2}(f(y))+\degtr_{k(f(y))}k(y),\qquad(y\in Y_1).
\end{equation}
Let us show that the functions $D_Y$ so defined satisfy conditions (i)--(iv) of 1.4. By (1.5.3), condition (ii) is the only nontrivial one. By (iv), one may clearly suppose that $Y=\overline{\{x\}}$. Then factor the structural morphism $s_Y$ of $Y$ through its image $s_Y(Y)=Z\hookrightarrow S$, say $s_Y=g\circ f$:
\[
\begin{tikzcd}
Y \arrow[r,"f"] & Z \arrow[r,hook,"g"] & S,
\end{tikzcd}
\]
where $Y$ and $Z$ are irreducible, $f$ is proper and surjective, and $g$ is a closed immersion. Put $x'=s_Y(x)$, $y'=s_Y(y)$. Then (1.5.2) and (1.5.1) give:
\[
D_Y(x)-D_Y(y)=\dimdot\underline{O}_{S,y'}-\dimdot\underline{O}_{S,x'}+\degtr_{k(x')}k(x)-\degtr_{k(y')}k(y).
\]
Since $S$ is catenary, one has

